\documentclass[11pt]{article}
\usepackage{fontspec}
\usepackage[margin=1in,top=1.1in,bottom=1.1in]{geometry}
\usepackage{booktabs}
\usepackage{longtable}
\usepackage{array}
\usepackage{amsmath}
\usepackage{xcolor}
\usepackage{titlesec}
\usepackage{fancyhdr}
\usepackage{enumitem}
\usepackage{caption}
\usepackage{unicode-math}
\usepackage[hidelinks]{hyperref}
\usepackage{microtype}
\usepackage{tabularx}
\usepackage{colortbl}
\usepackage{needspace}
\usepackage{float}

\setmainfont{Times New Roman}
\setmonofont{Consolas}[Scale=0.86]
\setmathfont{Cambria Math}

\definecolor{dqblue}{HTML}{1F3864}
\definecolor{dqgrey}{HTML}{5A5A5A}
\definecolor{dqrule}{HTML}{B4B4B4}
\definecolor{rowtint}{HTML}{F2F2F2}
\definecolor{warnbg}{HTML}{FBF3E4}
\definecolor{winbg}{HTML}{EAF1E4}

\titleformat{\section}{\Large\bfseries\color{dqblue}}{\thesection}{0.7em}{}
\titleformat{\subsection}{\large\bfseries\color{dqblue}}{\thesubsection}{0.6em}{}
\titleformat{\subsubsection}{\normalsize\bfseries}{\thesubsubsection}{0.5em}{}
\titlespacing{\section}{0pt}{1.6ex plus .3ex}{0.8ex}
\titlespacing{\subsection}{0pt}{1.3ex plus .3ex}{0.6ex}

\setlength{\headheight}{21.3pt}
\addtolength{\topmargin}{-7.3pt}
\setlength{\emergencystretch}{3em}
\pagestyle{fancy}
\fancyhf{}
\fancyhead[L]{\footnotesize\color{dqgrey}DynQuant Phase 1 --- External Comparison Campaign}
\fancyhead[R]{\footnotesize\color{dqgrey}\thepage}
\renewcommand{\headrulewidth}{0.4pt}
\renewcommand{\headrule}{\hbox to\headwidth{\color{dqrule}\leaders\hrule height \headrulewidth\hfill}}

\captionsetup{font=small,labelfont={bf,color=dqblue}}
\setlist{itemsep=0.25ex,topsep=0.5ex,parsep=0pt}

\newcommand{\x}{\texorpdfstring{$\times$}{x}}
\newcommand{\pt}{\,pts}
\newcommand{\pct}{\,\%}
\newcommand{\code}[1]{\texttt{#1}}
\newcommand{\minus}{$-$}
\newcommand{\e}[2]{$#1\!\times\!10^{#2}$}
\newcommand{\best}[1]{\textbf{#1}}

% a boxed aside, used for the load-bearing cautions
\newsavebox{\asidebox}
\newenvironment{aside}%
{\par\needspace{4\baselineskip}\vspace{0.8ex}%
\begin{lrbox}{\asidebox}\begin{minipage}{\dimexpr\textwidth-18pt\relax}\small}%
{\end{minipage}\end{lrbox}%
\noindent\setlength{\fboxsep}{9pt}\colorbox{warnbg}{\usebox{\asidebox}}%
\vspace{0.6ex}\par}

\begin{document}

%% ============================================================ TITLE
\begin{titlepage}
\centering
\vspace*{1.0in}
{\color{dqblue}\rule{\textwidth}{1.2pt}}\\[0.8em]
{\Huge\bfseries DynQuant --- Phase 1\par}
\vspace{0.5em}
{\LARGE The External Comparison Campaign\par}
\vspace{0.35em}
{\large DynQuant against GPTQ, AWQ, RTN and bitsandbytes NF4\par}
\vspace{0.7em}
{\color{dqblue}\rule{\textwidth}{1.2pt}}\\[1.8em]
{\large 45 comparison arms, 2 determinism re-runs, 5 packed-path records\par}
\vspace{0.3em}
{\large Two models, two tasks, one harness\par}
\vspace{2.2em}

\begin{minipage}{0.86\textwidth}
\small
\textbf{Summary.} This is the settled report of the campaign that measured DynQuant against
the three quantization methods it actually competes with --- GPTQ, AWQ and round-to-nearest
from \code{llm-compressor}, plus bitsandbytes NF4 --- at matched byte budgets, on two models,
two tasks and one shared evaluation harness. It exists because the paper's published
comparison was against numbers copied from other papers, and because the phrase ``4-bit
baseline'' turns out not to mean four bits.
\\[0.7em]
The campaign produced one outright win, a broad tie, and one clear loss, and it is the loss
that mattered most: at a 3-bit budget GPTQ was simultaneously \textbf{3.1\pct{} smaller and
1.34 points better} on the primary panel, and the same ordering replicated on all three
panels. It also separated DynQuant's two components for the first time: of the 25.78-point
margin over byte-matched round-to-nearest at 3.25 bits, \textbf{22.62 points are the
allocation prior and 3.16 points are the training signal}. Above four bits the signal is
worth nothing measurable.
\\[0.7em]
Every accuracy comparison is a paired exact McNemar test on stored per-example correctness,
every byte figure is computed on one convention for all methods, every baseline arm is
verified to have been genuinely rounded, and the packed CUDA path is confirmed to reproduce
the simulated one to within five predictions out of 3{,}080. Five defects found during the
campaign are documented in \S\ref{sec:defects}, including three that would each have produced
a wrong published number.
\end{minipage}

\vfill
{\small\color{dqgrey}
Measured 2026-07-26 through 2026-07-30 on a single NVIDIA A100 80\,GB.\\
Every number here was read back off the per-arm run records at write time, with one exception
flagged where it appears. \S\ref{sec:limits} states what is \emph{not} established.\\[0.4em]
\textbf{Superseded in part.} The 3-bit result in finding 3 was overturned on the primary panel
by the Phase 2 rebuild; see the companion report and the footnote in \S\ref{sec:replication}.\par}
\vspace{0.6em}
\end{titlepage}

\tableofcontents
\newpage

%% ============================================================ 1
\section{Bottom line}

Seven findings, in descending order of how much they should change your mind about the method.

\begin{enumerate}[leftmargin=1.4em]
\item \textbf{At 2.42\x{} compression DynQuant is lossless and GPTQ is not.} At byte-identical
size (1.4512 GiB against 1.4514 GiB) DynQuant scores 89.71\pct{} to GPTQ's 88.97\pct{} ---
\best{$+0.73$, $p = 0.017$} --- and is statistically indistinguishable from bf16
($-0.04$, $p = 0.905$) where GPTQ is not ($-0.77$, $p = 0.012$). This is the campaign's single
clean win. It is one tier, one task, one model.

\item \textbf{At 3.8\x{} compression DynQuant ties everything.} Across three panels and four
baselines, 11 of 12 comparisons against DynQuant's 4.25-bit arm fail to separate; the twelfth
(RTN on Qwen/Banking77, $+1.14$) favours DynQuant. On the tied-embedding model the tie is
achieved at \textbf{1.73\x{} fewer bytes} than the baselines' shipped configuration. The tie is
the result; the size is the win.

\item \textbf{At 4.9\x{} compression GPTQ wins, on all three panels.} $-1.34$
(\e{4.2}{-4}) on Qwen/CaseHOLD, $-1.04$ ($p = 0.011$) on Qwen/Banking77, $-0.78$
(\e{7.1}{-3}) on Mistral/Banking77. At matched bytes GPTQ is simultaneously \textbf{3.1\pct{}
smaller and 1.34 points better}. This was the method's real weakness and it replicated across
model, task and scale.\footnote{Reversed on Qwen/CaseHOLD by the Phase 2 rebuild:
$+1.54$ at 4.5\pct{} fewer bytes, $p < 0.0001$. The other two panels have not been re-run.
See \S\ref{sec:replication} and the Phase 2 report.}

\item \textbf{DynQuant beats AWQ at 3 bits by 2.5--3.3 points on both Qwen tasks} ---
$+3.33$ (\e{7.5}{-14}) and $+2.47$ (\e{3.3}{-7}) --- at half the bytes, and by $+3.39$
(\e{1.5}{-13}) against AWQ's matched-byte arm. Its clearest advantage over a calibration-based
method, though on Mistral at matched bytes the two are a dead tie ($+0.03$, $p = 1.00$).

\item \textbf{The margin over naive rounding is 88\pct{} architectural prior and 12\pct{}
training signal.} At 3.25 bits DynQuant beats byte-matched RTN by 25.78 points. Holding the
allocator, the graph, the floors and the budget fixed and replacing only the signal scores with
a constant splits that into \textbf{22.62 points of allocation and 3.16 points of signal}
($p = $ \e{1.9}{-15}). Both halves are real and separated; the sizes are very different.

\item \textbf{Above about 4 bits the training signal is worth nothing measurable.} At 4.25 bits
the same ablation moves 12 of 187 widths and buys $+0.19$ points ($p = 0.15$). The operating
rule that follows: DynQuant earns its keep only when the budget is tight enough that the role
floors cannot all be paid for. Above that point it is an expensive round-to-nearest.

\item \textbf{The accuracies survive real kernels and the bytes are exact; the speed is not
there yet.} Run packed on CUDA, the two Mistral arms differ from their simulated counterparts by
\textbf{5 and 1 predictions out of 3{,}080} (Qwen: exact, $4468 = 4468$ of 5{,}314). On-disk
size matches the accounted figure \textbf{to the byte} and peak VRAM is \textbf{3.55\x{} /
4.55\x} below bf16. But decode is 10--15\pct{} slower than bf16 at batches 1--8 and 3.1--5.7\x{}
slower at batch 32, from per-launch overhead and a missing large-$M$ path --- not from the bit
allocation. No baseline was run packed, so this report contains no speed comparison against
GPTQ or AWQ.
\end{enumerate}

%% ============================================================ 2
\section{Setup}

\subsection{Models and tasks}

Each model is paired with the task it was originally validated on in this project, so no pairing
is a fresh choice made after seeing results.

\begin{center}\small
\begin{tabular}{@{}lll@{}}
\toprule
 & \textbf{Qwen3.5-2B-Base} & \textbf{Mistral-7B-Instruct-v0.3} \\
\midrule
parameters & 1{,}881{,}825{,}088 & 7{,}248{,}023{,}552 \\
bf16 size & 3.5052 GiB & 13.5005 GiB \\
\code{embed\_tokens} / \code{lm\_head} & \textbf{tied}, 509{,}108{,}032 params & \textbf{untied}, 268{,}701{,}696 params \\
 & \quad = \textbf{27.054\pct} of the model & \quad = \textbf{3.707\pct} \\
task & CaseHOLD (legal holding, 5-way) & Banking77 (intent, 77-way) \\
eval set & 5{,}314 examples & 3{,}080 examples \\
also run on & Banking77, 3{,}080 examples & --- \\
\bottomrule
\end{tabular}
\end{center}

Three panels result: \textbf{Qwen/CaseHOLD} (primary, 21 arms), \textbf{Qwen/Banking77} (task
transfer, 10 arms) and \textbf{Mistral/Banking77} (scale and untied embedding, 14 arms).

Every arm in a panel starts from the same fine-tuned checkpoint and is scored through the same
\code{common.run\_eval} --- same prompt, same exemplars, same decode settings, same scorer. The
comparison is shared by construction, not by having been copied carefully.

\subsection{Baselines}

GPTQ, AWQ and RTN come from \textbf{llm-compressor} (the vLLM project's quantization library).
NF4 comes from \textbf{bitsandbytes}.

\begin{center}\small
\begin{tabular}{@{}ll@{}}
\toprule
calibration & 256 rows from the task's own train split, \code{seq\_len} 1024 \\
grouping & \code{group\_size} 128, \code{targets=["Linear"]} \\
symmetry & symmetric for GPTQ and RTN, asymmetric for AWQ \\
GPTQ & \code{dampening\_frac=0.01} (llm-compressor's default, named explicitly) \\
\code{ignore} & \code{["lm\_head"]} for the shipped-convention arms, \code{[]} for the \code{\_head} arms \\
\bottomrule
\end{tabular}
\end{center}

Two properties of this harness are load-bearing and are stated rather than assumed.

\subsubsection{Every baseline arm is verified to be actually rounded}

\code{llm-compressor}'s \code{oneshot} fits scales but does \emph{not} write rounded weights back
for any recipe except GPTQ --- so an in-process evaluation scores the \emph{original} checkpoint
and reports a suspiciously good number. The harness materializes the rounding itself, counts how
many modules moved, and then runs a fixed-point check on a spread of modules: re-quantizing an
already-quantized weight must be a no-op, and a tensor that was never rounded fails that. An arm
that fails raises \code{SystemExit} rather than reporting a number.

\code{weights\_moved} is recorded with every arm and is expected to be zero for GPTQ and equal to
the module count for RTN and AWQ. \textbf{Every arm matches that shape exactly:} across all 24
llm-compressor arms in the three panels, GPTQ reports \code{0/186}, \code{0/187} or \code{0/224}
with \code{max\_weight\_delta = 0.0} (already rounded), and every RTN and AWQ arm reports
\code{186/186}, \code{187/187} or \code{224/224} (rounded by the harness). \textbf{No arm scored
an unquantized checkpoint.}

\subsubsection{Bytes are accounted on one convention for all methods}

\code{accounted\_bytes} counts $\mathrm{numel} \times \mathrm{bits} + (\mathrm{numel} /
\mathrm{group\_size}) \times \mathrm{meta\_bits}$ for every quantized tensor and 16 bits per
parameter for everything left alone, including the embedding. This is the number quoted
throughout; it is not each library's self-report.

\subsection{How accuracy is compared}

Paired exact McNemar on per-example correctness, with a 95\pct{} confidence interval on the
paired difference. \code{flips a/b} counts examples the first arm got right and the second
wrong, and vice versa.

\begin{aside}
Two arms can have identical accuracy and still differ. On Qwen/CaseHOLD, GPTQ 4-bit and bf16
both score 89.74\pct{} with \textbf{58/58 flips} --- equal accuracy, 116 different predictions.
\textbf{Accuracy ties are not behavioural ties}, and nothing in this report claims otherwise.
\end{aside}

%% ============================================================ 3
\section{The accounting problem that has to be settled first}
\label{sec:accounting}

\begin{aside}
\textbf{A baseline labelled ``4-bit'' is not 4 bits, and on a tied-embedding model it is not
close.}
\end{aside}

\code{llm-compressor} targets \code{Linear} modules. \code{embed\_tokens} is an \code{Embedding},
so it is never touched; \code{lm\_head} is excluded by the default \code{ignore=["lm\_head"]}. On
Qwen3.5-2B those are the \emph{same tensor} --- $[248320, 2048]$, \textbf{508{,}559{,}360
parameters, 27.025\pct} of the model's 1{,}881{,}825{,}088 --- and it stays at 16 bits. Adding
the 1-D norm weights (106{,}304) brings the directly counted unquantized residue to
508{,}665{,}664, or \textbf{27.030\pct}.

The $f$ in the table below is the residue \emph{back-solved from the measured effective bits}
rather than counted, which is why it reads 27.054\pct:
$f = (7.3605 - 4.15625) / (16 - 4.15625)$. The two differ by about 442k parameters --- further
small modules \code{llm-compressor} declines to quantize --- and the gap is immaterial to every
conclusion here, but the counted and the fitted figure are not the same number and this report
does not imply they are.

The effective-bit arithmetic recovers this exactly. With \code{group\_size} 128, a 16-bit scale
and a 4-bit zero point, a quantized parameter costs $b + (16+4)/128 = b + 0.156$ bits. With an
fp16 residue fraction $f$:

\begin{center}\small
\begin{tabular}{@{}lrrrrr@{}}
\toprule
model & $f$ measured & $b{=}4$ predicted & $b{=}4$ measured & $b{=}3$ predicted & $b{=}3$ measured \\
\midrule
Qwen3.5-2B (tied) & 27.054\pct & 7.360 & \best{7.3605} & 6.625 & \best{6.6253} \\
Mistral-7B (untied) & 3.707\pct & 4.595 & \best{4.5953} & 3.625 & \best{3.6249} \\
\bottomrule
\end{tabular}
\end{center}

The \code{\_head} arms (\code{ignore=[]}) drive the residue to \textbf{0.029\pct} --- the
LayerNorm weights, which no method quantizes.

So the phrase ``4-bit GPTQ'' means 7.36 bits on the small tied model and 4.60 bits on the large
untied one. DynQuant quantizes every quantizable tensor including the embedding, so its 4.25-bit
map really stores 4.2486 bits.

\textbf{This is the single largest confound in the whole comparison and it is
architecture-dependent.} Every byte-ratio claim below therefore comes in two versions:

\begin{itemize}[leftmargin=1.4em]
\item \textbf{shipped convention} --- what a user actually gets from \code{llm-compressor} today,
\code{ignore=["lm\_head"]}.
\item \textbf{matched bytes} --- the \code{\_head} arms, \code{ignore=[]}, where the baselines
quantize the tie too and DynQuant's byte advantage disappears by construction.
\end{itemize}

Both are reported. Only the matched-byte comparison answers \emph{is the allocation better};
only the shipped-convention one answers \emph{what happens if I use these tools as they ship}.

%% ============================================================ 4
\section{Panel A --- Qwen3.5-2B-Base / CaseHOLD (primary, 21 arms)}

bf16 reference: \textbf{89.74\pct} (4769/5314), 3.505 GiB.

\begin{center}\footnotesize
\begin{tabular}{@{}llrrrrrrr@{}}
\toprule
arm & method & eff. bits & GiB & ratio & acc \% & correct & $\Delta$ bf16 & $p$ \\
\midrule
gptq\_4b\_head & GPTQ, tie quantized & 4.1597 & 0.9113 & 3.85\x & \best{89.76} & 4770 & $+0.02$ & 1.00 \\
gptq\_4b & GPTQ, shipped & 7.3605 & 1.6125 & 2.17\x & 89.74 & 4769 & $+0.00$ & 1.00 \\
fp16 & --- & 16 & 3.505 & 1.00\x & 89.74 & 4769 & --- & --- \\
\rowcolor{winbg}
\textbf{dq\_iso6p63} & \textbf{DynQuant} & 6.6247 & 1.4512 & 2.42\x & \best{89.71} & 4767 & $-0.04$ & 0.905 \\
\textbf{dq\_iso7p36} & \textbf{DynQuant} & 7.3602 & 1.6123 & 2.17\x & 89.63 & 4763 & $-0.11$ & 0.327 \\
\textbf{dq\_4p25} & \textbf{DynQuant} & 4.2486 & 0.9307 & 3.77\x & \best{89.25} & 4743 & $-0.49$ & 0.099 \\
awq\_4b & AWQ, shipped & 7.3605 & 1.6125 & 2.17\x & 89.18 & 4739 & $-0.56$ & 0.027 \\
dq\_ctl4p25 & \emph{control} (no signal) & 4.2486 & 0.9307 & 3.77\x & 89.07 & 4733 & $-0.68$ & 0.021 \\
nf4\_4b & bnb NF4 & 7.3391 & 1.6078 & 2.18\x & 89.01 & 4730 & $-0.73$ & \e{4.2}{-3} \\
gptq\_3b & GPTQ, shipped & 6.6253 & 1.4514 & 2.41\x & 88.97 & 4728 & $-0.77$ & 0.012 \\
rtn\_4b & RTN, shipped & 7.3605 & 1.6125 & 2.17\x & 88.93 & 4726 & $-0.81$ & \e{5.6}{-3} \\
awq\_4b\_head & AWQ, tie quantized & 4.1597 & 0.9113 & 3.85\x & 88.93 & 4726 & $-0.81$ & \e{1.2}{-3} \\
rtn\_4b\_head & RTN, tie quantized & 4.1597 & 0.9113 & 3.85\x & 88.80 & 4719 & $-0.94$ & \e{1.8}{-3} \\
gptq\_3b\_head & GPTQ, tie quantized & 3.1522 & 0.6906 & 5.08\x & 88.03 & 4678 & $-1.71$ & \e{1.1}{-7} \\
\textbf{dq\_3p25} & \textbf{DynQuant} & 3.2494 & 0.7118 & 4.92\x & \best{86.70} & 4607 & $-3.05$ & \e{4.9}{-17} \\
dq\_ctl3p25 & \emph{control} (no signal) & 3.2494 & 0.7118 & 4.92\x & 83.53 & 4439 & $-6.21$ & \e{1.2}{-44} \\
awq\_3b & AWQ, shipped & 6.6253 & 1.4514 & 2.41\x & 83.36 & 4430 & $-6.38$ & \e{3.7}{-44} \\
awq\_3b\_head & AWQ, tie quantized & 3.1522 & 0.6906 & 5.08\x & 83.31 & 4427 & $-6.44$ & \e{1.6}{-45} \\
dq\_ctl3p25\_emb3 & \emph{control} + 3b embed pin & 3.2494 & 0.7118 & 4.92\x & 80.56 & 4281 & $-9.18$ & \e{1.1}{-73} \\
rtn\_3b & RTN, shipped & 6.6253 & 1.4514 & 2.41\x & 65.53 & 3482 & $-24.22$ & \e{2.1}{-259} \\
rtn\_3b\_head & RTN, tie quantized & 3.1522 & 0.6906 & 5.08\x & 60.91 & 3237 & $-28.83$ & ${\sim}0$ \\
\bottomrule
\end{tabular}
\end{center}

Read the ordering with care: \textbf{it is not a size ordering}. \code{dq\_3p25} at 0.7118 GiB
sits between arms that are twice its size in both directions. The size column is what makes this
table a comparison rather than a leaderboard.

\subsection*{Three observations}

\textbf{The 4-bit tier is flat.} Not one adjacent gap between \code{gptq\_4b\_head} (89.76) and
\code{rtn\_4b\_head} (88.80) separates statistically. Six methods spanning rounding-only to
inverse-Hessian error feedback land inside one point of each other. Whatever distinguishes these
algorithms is not visible at 4 bits on this task.

\textbf{The 3-bit tier is where methods separate, and it separates by a lot.} From 88.03 (GPTQ)
to 60.91 (RTN) is 27 points at identical bytes. This is the regime the comparison is actually
about.

\textbf{RTN's collapse is the scale of the prize.} RTN at 3 bits loses 24--29 points. Everything
any of these methods does --- Hessian error feedback, activation smoothing, signal-driven
allocation --- is competing for that 27-point gap.

%% ============================================================ 5
\section{Panel B --- Qwen3.5-2B-Base / Banking77 (task transfer, 10 arms)}

Same model, same allocator, same fine-tuning recipe, different task. bf16: \textbf{93.41\pct}
(2877/3080).

\begin{center}\small
\begin{tabular}{@{}lrrrrrr@{}}
\toprule
arm & eff. bits & GiB & acc \% & correct & $\Delta$ bf16 & $p$ \\
\midrule
fp16 & 16 & 3.505 & 93.41 & 2877 & --- & --- \\
gptq\_4b & 7.3605 & 1.6125 & 93.08 & 2867 & $-0.32$ & 0.076 \\
nf4\_4b & 7.3391 & 1.6078 & 92.89 & 2861 & $-0.52$ & 0.037 \\
\textbf{dq\_4p25} & 4.2486 & 0.9307 & \best{92.66} & 2854 & $-0.75$ & \e{2.7}{-3} \\
awq\_4b & 7.3605 & 1.6125 & 92.50 & 2849 & $-0.91$ & \e{2.3}{-4} \\
gptq\_3b & 6.6253 & 1.4514 & 92.14 & 2838 & $-1.27$ & \e{2.2}{-5} \\
rtn\_4b & 7.3605 & 1.6125 & 91.53 & 2819 & $-1.88$ & \e{4.3}{-9} \\
\textbf{dq\_3p25} & 3.2494 & 0.7118 & \best{91.10} & 2806 & $-2.31$ & \e{6.8}{-10} \\
awq\_3b & 6.6253 & 1.4514 & 88.64 & 2730 & $-4.77$ & \e{2.0}{-25} \\
rtn\_3b & 6.6253 & 1.4514 & 61.62 & 1898 & $-31.79$ & \e{2.7}{-255} \\
\bottomrule
\end{tabular}
\end{center}

The shape of Panel A survives: the 4-bit tier is tight, RTN collapses at 3 bits, DynQuant's
4.25-bit arm sits inside the 4-bit cluster at 1.73\x{} fewer bytes, and its 3.25-bit arm sits
between GPTQ and AWQ. The \code{\_head} arms were not re-run here --- the accounting lesson from
Panel A carries over unchanged and re-measuring it would have bought nothing.

%% ============================================================ 6
\section{Panel C --- Mistral-7B-Instruct-v0.3 / Banking77 (scale, untied, 14 arms)}

bf16: \textbf{94.38\pct} (2907/3080), 13.50 GiB.

\begin{center}\small
\begin{tabular}{@{}lrrrrrrr@{}}
\toprule
arm & eff. bits & GiB & ratio & acc \% & correct & $\Delta$ bf16 & $p$ \\
\midrule
\textbf{dq\_iso4p60} & 4.5949 & 3.8770 & 3.48\x & \best{94.42} & 2908 & $+0.03$ & 1.00 \\
awq\_4b & 4.5953 & 3.8774 & 3.48\x & \best{94.42} & 2908 & $+0.03$ & 1.00 \\
gptq\_3b & 3.6249 & 3.0586 & 4.41\x & 94.38 & 2907 & $+0.00$ & 1.00 \\
fp16 & 16 & 13.50 & 1.00\x & 94.38 & 2907 & --- & --- \\
rtn\_4b & 4.5953 & 3.8774 & 3.48\x & 94.25 & 2903 & $-0.13$ & 0.503 \\
nf4\_4b & 4.5671 & 3.8536 & 3.50\x & 94.25 & 2903 & $-0.13$ & 0.481 \\
gptq\_4b & 4.5953 & 3.8774 & 3.48\x & 94.25 & 2903 & $-0.13$ & 0.424 \\
\textbf{dq\_4p25} & 4.2500 & 3.5859 & 3.76\x & \best{94.19} & 2901 & $-0.19$ & 0.263 \\
dq\_4p25\_ft1map & 4.2500 & 3.5859 & 3.76\x & 94.19 & 2901 & $-0.19$ & 0.263 \\
\textbf{dq\_iso3p62} & 3.6244 & 3.0581 & 4.41\x & 94.06 & 2897 & $-0.32$ & 0.154 \\
awq\_3b & 3.6249 & 3.0586 & 4.41\x & 94.03 & 2896 & $-0.36$ & 0.090 \\
\textbf{dq\_3p25} & 3.2500 & 2.7422 & 4.92\x & \best{93.60} & 2883 & $-0.78$ & \e{3.2}{-3} \\
dq\_3p25\_ft1map & 3.2500 & 2.7422 & 4.92\x & 93.57 & 2882 & $-0.81$ & \e{2.6}{-3} \\
rtn\_3b & 3.6249 & 3.0586 & 4.41\x & 93.34 & 2875 & $-1.04$ & \e{5.8}{-5} \\
\bottomrule
\end{tabular}
\end{center}

\textbf{This panel has little discriminating power, and that is its finding.} The whole spread
from best to worst is \textbf{1.08 points} on 3{,}080 examples. Even RTN at 3 bits --- which
loses 24 and 32 points on the two Qwen panels --- loses only 1.04 here. Two explanations, and the
panel cannot separate them:

\begin{itemize}[leftmargin=1.4em]
\item \textbf{Scale.} 7.25\,B parameters have redundancy that 1.88\,B do not, so the same
relative weight perturbation costs less accuracy.
\item \textbf{The untied embedding.} The tensor that dominates the tied model's error budget is
3.7\pct{} of this one, and the baselines leave it in fp16 either way.
\end{itemize}

The panel is not \emph{entirely} flat, and the distinction matters. Six comparisons do separate
here, all in the 3-bit tier: \code{rtn\_3b} vs bf16 ($-1.04$, \e{5.8}{-5}), \code{dq\_3p25} vs
bf16 ($-0.78$, \e{3.2}{-3}), \code{dq\_3p25} vs \code{gptq\_3b} ($-0.78$, \e{7.1}{-3}),
\code{dq\_iso3p62} vs \code{rtn\_3b} ($+0.71$, \e{7.2}{-3}), \code{dq\_3p25} vs \code{dq\_4p25}
($-0.58$, $p = 0.027$) and \code{dq\_iso3p62} vs \code{dq\_3p25} ($+0.45$, $p = 0.039$). So the
panel has enough power to resolve a 0.5-point effect; what it lacks is \textbf{dynamic range}.
Every 4-bit-class comparison fails to separate, and a 25-point effect and a 0.7-point effect
cannot look different in kind when the ceiling is one point away. The panel can rank the 3-bit
arms; it cannot tell you how much the ranking is worth.

The practical consequence is that a 7B model on a 77-way classification task is a poor
discriminating benchmark for weight quantization. Any method that rounds competently passes.
\textbf{A comparison that had only been run here would have concluded that all five methods are
equivalent.}

Note also that the byte picture inverts relative to Panel A: DynQuant's advantage over the
shipped convention drops from 1.73\x{} to \textbf{1.08\x}, because the fp16 residue the baselines
carry drops from 27.0\pct{} to 3.7\pct. DynQuant's storage advantage is largely an artifact of
\emph{what the baselines decline to quantize}, and that is an architectural property, not a
property of either method.

%% ============================================================ 7
\section{Matched bytes --- the six budgets that answer the actual question}

Every comparison below is between arms of the same size, so allocation is the only difference
left.

\subsection{1.61 GiB / 7.36 bits / 2.17\x{} (Qwen/CaseHOLD)}

\begin{center}\small
\begin{tabular}{@{}lrlrr@{}}
\toprule
pair & delta & CI95 & flips & $p$ \\
\midrule
dq\_iso7p36 vs gptq\_4b & $-0.11$ & $[-0.51, +0.28]$ & 54/60 & 0.64 \\
dq\_iso7p36 vs awq\_4b & $+0.45$ & $[-0.02, +0.92]$ & 93/69 & 0.070 \\
dq\_iso7p36 vs rtn\_4b & \best{$+0.70$} & $[+0.13, +1.26]$ & 137/100 & \best{0.019} \\
dq\_iso7p36 vs nf4\_4b & \best{$+0.62$} & $[+0.13, +1.11]$ & 105/72 & \best{0.016} \\
dq\_iso7p36 vs fp16 & $-0.11$ & $[-0.30, +0.08]$ & 10/16 & 0.327 \\
\bottomrule
\end{tabular}
\end{center}

DynQuant ties GPTQ, edges AWQ without separating, and beats both RTN and NF4. Lossless relative
to bf16.

\subsection{1.45 GiB / 6.63 bits / 2.42\x{} (Qwen/CaseHOLD) --- the one outright win}
\label{sec:663}

\begin{center}\small
\begin{tabular}{@{}lrlrr@{}}
\toprule
pair & delta & CI95 & flips & $p$ \\
\midrule
\rowcolor{winbg}
\textbf{dq\_iso6p63 vs gptq\_3b} & \best{$+0.73$} & $[+0.14, +1.32]$ & 147/108 & \best{0.017} \\
dq\_iso6p63 vs awq\_3b & $+6.34$ & $[+5.43, +7.25]$ & 474/137 & \e{2.1}{-44} \\
dq\_iso6p63 vs rtn\_3b & $+24.18$ & $[+22.71, +25.65]$ & 1434/149 & \e{6.9}{-264} \\
dq\_iso6p63 vs fp16 & $-0.04$ & $[-0.35, +0.27]$ & 34/36 & 0.905 \\
\bottomrule
\end{tabular}
\end{center}

\textbf{This is the result the method should be presented on.} At 1.4512 GiB DynQuant is
indistinguishable from bf16; GPTQ at 1.4514 GiB is 0.77 points below it ($p = 0.012$). Same
bytes, same eval, and one of the two is lossless.

The honest qualification: 6.63 bits is a mild budget. \code{dq\_iso6p63} and \code{dq\_iso7p36}
do not differ from each other ($+0.08$, $p = 0.72$), so the win is not ``DynQuant degrades
gracefully'' --- it is ``DynQuant has not started degrading yet at a budget where GPTQ has.''

\subsection{0.91--0.93 GiB / 4.16--4.25 bits / 3.8\x{} (Qwen/CaseHOLD)}

\begin{center}\small
\begin{tabular}{@{}lrlrr@{}}
\toprule
pair & delta & CI95 & flips & $p$ \\
\midrule
dq\_4p25 vs gptq\_4b\_head & $-0.51$ & $[-1.08, +0.06]$ & 106/133 & 0.092 \\
dq\_4p25 vs awq\_4b\_head & $+0.32$ & $[-0.25, +0.89]$ & 127/110 & 0.299 \\
dq\_4p25 vs rtn\_4b\_head & $+0.45$ & $[-0.16, +1.06]$ & 150/126 & 0.166 \\
\bottomrule
\end{tabular}
\end{center}

Nothing separates. DynQuant is also \textbf{2.1\pct{} larger} than the GPTQ arm here, so on bytes
it loses slightly.

\subsection{0.69--0.71 GiB / 3.15--3.25 bits / 5.0\x{} (Qwen/CaseHOLD) --- the weakness}
\label{sec:weakness}

\begin{center}\small
\begin{tabular}{@{}lrlrr@{}}
\toprule
pair & delta & CI95 & flips & $p$ \\
\midrule
\textbf{dq\_3p25 vs gptq\_3b\_head} & \best{$-1.34$} & $[-2.07, -0.60]$ & 162/233 & \best{\e{4.2}{-4}} \\
dq\_3p25 vs awq\_3b\_head & $+3.39$ & $[+2.49, +4.29]$ & 388/208 & \e{1.5}{-13} \\
dq\_3p25 vs rtn\_3b\_head & $+25.78$ & $[+24.26, +27.30]$ & 1531/161 & \e{3.0}{-280} \\
\bottomrule
\end{tabular}
\end{center}

\textbf{GPTQ is 3.1\pct{} smaller and 1.34 points better.} There is no reading of this budget in
which DynQuant wins it. The mechanism is not mysterious: GPTQ propagates and compensates
quantization error through an inverse-Hessian update, and at the point where every module is
badly quantized, compensating the error beats choosing which modules to damage. DynQuant chooses
well and then rounds naively; GPTQ chooses nothing and then rounds cleverly. At 3 bits the second
matters more.

\begin{aside}
\textbf{This budget is the whole of Phase 2.} The companion report describes the rebuild that
reverses this row on this panel --- $+1.54$ at 4.5\pct{} fewer bytes, $p < 0.0001$ --- without
adopting error feedback, an inverse Hessian, or a calibration set.
\end{aside}

\subsection{3.88 GiB and 3.06 GiB (Mistral/Banking77)}

Two more matched-byte budgets on the untied 7B model. \code{dq\_iso4p60} sits at 3.8770 GiB
against the baselines' 3.8774; \code{dq\_iso3p62} at 3.0581 against 3.0586.

\begin{center}\small
\begin{tabular}{@{}lrlrr@{}}
\toprule
pair & delta & CI95 & flips & $p$ \\
\midrule
dq\_iso4p60 vs gptq\_4b & $+0.16$ & $[-0.12, +0.44]$ & 12/7 & 0.36 \\
dq\_iso4p60 vs awq\_4b & $+0.00$ & $[-0.24, +0.24]$ & 7/7 & 1.00 \\
dq\_iso4p60 vs rtn\_4b & $+0.16$ & $[-0.10, +0.42]$ & 11/6 & 0.33 \\
dq\_iso4p60 vs nf4\_4b & $+0.16$ & $[-0.10, +0.42]$ & 11/6 & 0.33 \\
dq\_iso4p60 vs bf16 & $+0.03$ & $[-0.20, +0.26]$ & 7/6 & 1.00 \\
\midrule
dq\_iso3p62 vs gptq\_3b & $-0.32$ & $[-0.81, +0.16]$ & 24/34 & 0.24 \\
dq\_iso3p62 vs awq\_3b & $+0.03$ & $[-0.45, +0.51]$ & 29/28 & 1.00 \\
\textbf{dq\_iso3p62 vs rtn\_3b} & \best{$+0.71$} & $[+0.21, +1.22]$ & 42/20 & \best{\e{7.2}{-3}} \\
dq\_iso3p62 vs bf16 & $-0.32$ & $[-0.73, +0.08]$ & 15/25 & 0.15 \\
\bottomrule
\end{tabular}
\end{center}

At 4.60 bits \code{dq\_iso4p60} is the joint-best arm in the panel (94.42\pct, tied with AWQ at
2908/3080 and indistinguishable from bf16) --- which on a 1.08-point spread means only that it is
not worse.

\textbf{The 3.62-bit row is the one that matters, and it changes a verdict.} At matched bytes
DynQuant beats RTN by \best{$+0.71$} (\e{7.2}{-3}) on Mistral. The shipped-convention comparison
had \emph{failed} to separate ($+0.26$, $p = 0.42$) --- but that pitted a 3.25-bit DynQuant
against a 3.62-bit RTN, giving away 10\pct{} of the size. Match the bytes and the separation
appears. It also confirms GPTQ's edge is not scale-specific: $-0.32$ here, in the same direction
as both Qwen panels, though this panel is too compressed to separate it.

\subsection{Summary across all four Qwen budgets}

\begin{center}\small
\begin{tabular}{@{}lrrrrrrl@{}}
\toprule
budget & ratio & DynQuant & GPTQ & AWQ & RTN & NF4 & verdict \\
\midrule
7.36 b & 2.17\x & 89.63 & 89.74 & 89.18 & 88.93 & 89.01 & tie GPTQ, beat RTN/NF4 \\
\rowcolor{winbg}
6.63 b & 2.42\x & \best{89.71} & 88.97 & 83.36 & 65.53 & --- & \textbf{DynQuant wins} \\
4.16 b & 3.85\x & 89.25 & 89.76 & 88.93 & 88.80 & --- & tie \\
3.15 b & 5.08\x & 86.70 & \best{88.03} & 83.31 & 60.91 & --- & \textbf{GPTQ wins} \\
\bottomrule
\end{tabular}
\end{center}

\textbf{DynQuant's advantage is a band, not a trend.} It is absent at mild budgets, present at
6.63 bits, absent again at 4.16, and negative at 3.15.

%% ============================================================ 8
\section{Do the accounted bytes exist? The packed path, measured}
\label{sec:packed}

Every accuracy number in \S4--\S7 comes from the \textbf{simulated} path: the quantizer writes
dequantized values back in place, so the arithmetic is quantized but the tensors are still bf16.
Every byte number is \emph{accounted} --- computed from the width map. Both need checking against
a device, and on Mistral they have been.

\subsection{Provenance first, because these records predate the corrected bit map}

They were written 2026-07-29 03:21--04:27; the map was refit 2026-07-30 11:52 (defect 2,
\S\ref{sec:defects}). Diffing the two maps module by module:

\begin{center}\small
\begin{tabular}{@{}llrl@{}}
\toprule
target & width histogram & \code{nbytes} & placement vs shipped map \\
\midrule
4.25 b & \code{\{3: 10, 4: 212, 8: 4\}} --- identical & 3{,}850{,}371{,}072 & \textbf{0 of 226 modules differ} \\
3.25 b & \code{\{2: 35, 3: 137, 4: 54\}} --- identical & 2{,}944{,}401{,}408 & 24 of 226 differ \\
\bottomrule
\end{tabular}
\end{center}

The 4.25-bit map is bit-identical to the shipped one. The 3.25-bit one has the same histogram and
the same byte total with 24 modules placed differently --- so its size and VRAM figures transfer
unchanged. More usefully, both packed arms are \textbf{map-matched to the \code{\_ft1map} arms in
Panel C}, which makes the kernel comparison below clean.

\subsection{The accounting is exact, not approximate}

\begin{center}\small
\begin{tabular}{@{}lrrrr@{}}
\toprule
arm & accounted \code{nbytes} & \code{packed\_bytes} on disk & resident VRAM & over manifest \\
\midrule
dense bf16 & 14{,}496{,}047{,}616 & 14{,}495{,}514{,}624 & 13.5005 GiB & --- \\
packed 4.25 & 3{,}850{,}371{,}072 & \best{3{,}850{,}371{,}072} & 3.5864 GiB & $+0.014\%$ \\
packed 3.25 & 2{,}944{,}401{,}408 & \best{2{,}944{,}401{,}408} & 2.7427 GiB & $+0.018\%$ \\
\bottomrule
\end{tabular}
\end{center}

Both packed arms match the accounted figure \textbf{to the byte} on disk, \code{manifest\_bytes}
agrees, and resident VRAM exceeds it by 0.0005 GiB --- the packing metadata and the fp16 norms
the format keeps unquantized. \S\ref{sec:accounting}'s arithmetic is not a model of the
checkpoint; it is the checkpoint, and it holds on a device.

\subsection{The kernels reproduce the simulated path, which is what licenses \S4--\S6}

Comparing each packed arm against the simulated arm built from the same map, at the prediction
level rather than by net accuracy:

\begin{center}\small
\begin{tabular}{@{}lrrrr@{}}
\toprule
pair & packed & simulated & flips & $p$ \\
\midrule
packed 4.25 vs \code{dq\_4p25\_ft1map} & 2900/3080 & 2901/3080 & \textbf{2 / 3} & 1.00 \\
packed 3.25 vs \code{dq\_3p25\_ft1map} & 2881/3080 & 2882/3080 & \textbf{0 / 1} & 1.00 \\
\bottomrule
\end{tabular}
\end{center}

Five differing predictions out of 3{,}080 at 4.25 bits and \textbf{one} at 3.25 bits. On
Qwen/CaseHOLD the earlier campaign got exact agreement --- $4468 = 4468$ on 5{,}314 problems ---
alongside 417 kernel-parity tests and a clean four-tool \code{compute-sanitizer} sweep. So the
simulated accuracies this report is built on are what the real kernels produce, and that is
measured rather than assumed.

\subsection{Peak VRAM tracks the manifest}

This is the claim the original paper explicitly conceded it could not make.

\begin{center}\small
\begin{tabular}{@{}lrrrrrr@{}}
\toprule
arm & peak, batch 1 & vs bf16 & batch 4 & batch 8 & batch 32 & full 3080-row eval \\
\midrule
dense bf16 & 13.6766 GiB & --- & 14.1810 & 14.8536 & 18.8891 & 19.49 GiB $\dagger$ \\
packed 4.25 & 3.8526 GiB & \best{3.55\x} & 4.2993 & 4.9395 & 8.9750 & 9.5759 GiB \\
packed 3.25 & 3.0089 GiB & \best{4.55\x} & 3.4556 & 4.0958 & 8.1313 & 8.7322 GiB \\
\bottomrule
\end{tabular}
\end{center}

$\dagger$ The two packed eval peaks are read off the stage-6 records; the dense one is carried
from the Mistral write-up, because the dense record on the box holds \code{peak\_decode\_gib} but
no full-eval peak. \textbf{It is the only number in this report not read back off disk at write
time.}

At batch 1 the overhead above resident weights is \textbf{0.2662 GiB in both packed arms} --- the
same figure to four decimals, so it is the KV cache and activations for a 721-token prompt, not
per-arm slack. But the batch-32 and full-eval columns are the ones a serving decision turns on:
\textbf{3.55\x{} on weights is 2.04\x{} at eval peak}, because activations and KV cache are not
quantized. Quantization buys headroom in proportion to how much of your memory is weights, and
that fraction is not 100\pct.

\subsection{Decode is slower than bf16, and the batch-32 case is severe}

\begin{center}\small
\begin{tabular}{@{}lrrrr@{}}
\toprule
arm & batch 1 & batch 4 & batch 8 & batch 32 \\
\midrule
dense bf16 & 42.41 tok/s & 165.50 & 328.01 & \best{1244.41} \\
packed 4.25 & 36.90 (0.87\x) & 145.15 (0.88\x) & 294.03 (0.90\x) & 400.48 (\best{0.32\x}) \\
packed 3.25 & 36.93 (0.87\x) & 149.30 (0.90\x) & 280.01 (0.85\x) & 220.30 (\best{0.18\x}) \\
\bottomrule
\end{tabular}
\end{center}

Two cautions on the batch 1--8 rows. These are single runs per arm, measured about half an hour
apart, and the Qwen campaign established that batch-1 decode has a \textbf{12\pct{} spread across
repeated bf16 runs in one session} --- wide enough that a 0.87\x{} ratio from unrepeated cells is
not separable from host jitter. An earlier version of that Qwen table reported 0.98\x{} from
exactly this mistake and the repeated measurement made it 0.90\x. So read 0.85--0.90\x{} as
``slower, magnitude not pinned.''

What does not depend on the bf16 baseline is the comparison \emph{between} the two packed arms,
and it is the informative one: at batch 1 they decode in \textbf{27.099 ms and 27.078 ms ---
0.08\pct{} apart --- while carrying 3.5864 and 2.7427 GiB of weights}. A quarter less weight
traffic produced no measurable change. This kernel is not bandwidth-bound; packed moves 3.85 GB
in 27.1 ms = 142 GB/s against roughly 1935 GB/s of A100 HBM. The deficit at batches 1--8 is a
\emph{fixed additive} 2.6--4.2 ms/step, flat while throughput varies eightfold --- per-launch
cost, about 15 \textmu s across 226 quantized modules, which is what CUDA Graphs exists to
remove.

Batch 32 is a different failure with a known cause, not a worse version of the same one:
\code{GEMV\_MAX\_ROWS} is \textbf{8}, so batches 1--8 use the packed GEMV and 32 falls back to
dequantize-then-\code{F.linear}, materialising a bf16 copy per call and therefore reading
\emph{more} memory than the dense arm. 25.7 ms/step becomes 79.9 ms and 145.3 ms. That is a
missing tensor-core GEMM path, and it is why the narrower arm is the slower one there.

Prefill is unaffected: at batch 32 (23{,}072 tokens) it is 1.7924 s dense against 1.8013 s and
1.8573 s packed, within 4\pct, because prefill already goes through dequantize-then-GEMM. Packing
cost \textbf{1685 s and 1746 s} --- 28--29 minutes for a 7\,B model against the plan's gate of
under 6 minutes for a 14\,B one, roughly an order of magnitude off per parameter.

\begin{aside}
\textbf{The summary.} The size and VRAM claims are real, exact, and confirmed on hardware, and
the accuracy claims survive the move to real kernels. \textbf{The speed claim does not exist
yet} --- decode is slower at every batch size tested, and the reason is per-launch overhead and a
missing large-$M$ path, neither of which is a property of the bit allocation this report
evaluates.
\end{aside}

%% ============================================================ 9
\section{What is the training signal actually worth?}
\label{sec:ablation}

DynQuant has two components that could explain its results: an \textbf{architectural prior}
(role-based floors, structural protection of the recurrence coefficients, honest byte accounting)
and a \textbf{training-time signal} (activation saliency $\times$ gradient plasticity, ranked and
fed to the allocator). These are separable, and separating them is the most informative thing in
the campaign.

\subsection{The control}

The allocator runs the knapsack twice at every target: once with the real scores, once with
\code{uniform = dict.fromkeys(scores, 0.5)} --- graph, floors, group size and budget held fixed.
\textbf{This is the only comparison in the work where exactly one input changes.}

It is not a straw man. With equal scores the ROI metric $\mathrm{score} / (\mathrm{params} \times
\Delta\mathrm{bits})$ degenerates to pure size, so the control is a real strategy: \emph{spend
bits where they are cheapest per parameter}. All arms are byte-exact --- net $\mathrm{params}
\times \mathrm{bits}$ between the maps is exactly $+0$, because the budget binds. Different width
histograms at identical size is expected, not a bug.

\subsection{The result}

\begin{center}\small
\begin{tabular}{@{}lrrr@{}}
\toprule
target & widths the signal moves & signal worth & $p$ \\
\midrule
4.25 b & 12 / 187 (6\pct) & $+0.19$ & 0.15 \\
3.25 b & 87 / 187 (47\pct) & \best{$+3.16$} & \best{\e{1.9}{-15}} \\
\bottomrule
\end{tabular}
\end{center}

At 3.25 bits the full decomposition, all three arms at exactly 0.7118 GiB:

\begin{center}\small
\begin{tabular}{@{}lrl@{}}
\toprule
arm & acc \% & what the step buys \\
\midrule
RTN 3-bit + head & 60.91 & \\
 & \quad $\downarrow\ +22.62$ & role floors, structural protection, ROI knapsack (\e{2.8}{-229}) \\
uniform-score control & 83.53 & \\
 & \quad $\downarrow\ +3.16$ & the training signal (\e{1.9}{-15}) \\
DynQuant & 86.70 & \\
\bottomrule
\end{tabular}
\end{center}

\textbf{About 88\pct{} architectural prior, 12\pct{} training signal.} Both halves separate
overwhelmingly. The sharpest way to say it: \emph{without the signal, DynQuant at 3.25 bits is
AWQ} --- 83.53\pct{} against \code{awq\_3b\_head}'s 83.31\pct{} ($p = 0.66$), and AWQ's arm is
3\pct{} smaller.

At 4.25 bits the signal is worth nothing measurable, and the reason is mechanical rather than
mysterious: at that budget the role floors are nearly all affordable, so the allocator has almost
no decisions left to make. \textbf{The signal's value tracks how many decisions it is allowed to
make.}

\begin{aside}
\textbf{Operating rule.} Run DynQuant when the target is tight enough that the role floors cannot
all be satisfied. Above that point it is an expensive round-to-nearest.\\[0.4em]
This rule is also the reason Phase 2 targeted the 3-bit tier and nothing else, and the reason its
central finding --- that allocation granularity multiplies the signal's value --- is a direct
extension of this measurement: at \emph{row} granularity the same signal is handed 571{,}968
decisions instead of 187.
\end{aside}

\subsection{The decision that looked like it should carry the 3.16, and didn't}

The control breaches the tied embedding's inherited 8-bit floor all the way down to 2 bits; the
real map stops at 3. That one difference is 27\pct{} of the model and was the obvious candidate
for the whole margin. It was tested directly by pinning the control's embedding to 3 bits at the
same budget:

\begin{center}\small
\begin{tabular}{@{}lrlrr@{}}
\toprule
3.25 b, all at 0.7118 GiB & embedding & other 186 modules & acc \% & correct \\
\midrule
DynQuant (signal) & 3 b & signal-ROI & \best{86.70} & 4607/5314 \\
uniform-score control & 2 b & size-ROI & 83.53 & 4439 \\
uniform-score control + 3-bit pin & 3 b & size-ROI & \best{80.56} & 4281 \\
\bottomrule
\end{tabular}
\end{center}

\textbf{It landed below both.} Pinning the embedding to 3 bits made the control 2.97 points
\emph{worse} --- CI $[-3.78, -2.17]$, flips 100/258, $p = $ \e{2.4}{-12} --- and 6.13 points
behind the real map ($p = $ \e{4.5}{-45}).

The rescue has to be paid for, and \emph{how} it is paid dominates it. Under size-ROI, buying the
embedding's third bit spreads 55 modules from 4 to 3 bits, and that loses more than the embedding
gains. Under signal-ROI the same purchase is financed by holding 32 modules at 4 bits and driving
24 down to 2 --- and that earns it back. \textbf{What the signal supplies at this budget is the
payment schedule, not the embedding decision.}

Two caveats, both real:
\begin{itemize}[leftmargin=1.4em]
\item \textbf{It measures a pair, not the pin.} Changing the embedding's width at a fixed budget
necessarily changes other widths. The arm measures (embedding 3 b $+$ its financing), not the
embedding bit.
\item \textbf{The $2\times2$ has three cells, not four.} The fourth --- signal scores with the
embedding \emph{forced} to 2 --- cannot be built: \code{AllocationPolicy} exposes floors and
\code{structural\_roles}, never a per-module ceiling.
\end{itemize}

\subsection{An untested lever this exposed}

Lowering the embedding floor from 8 to 2 under \emph{real} scores produces
\code{\{2:8, 3:65, 4:78, 8:36\}} at 3.2489 average bits, where the downgrade-from-floor path
produces \code{\{2:25, 3:94, 4:32, 8:36\}} at the same budget. Same scores, same budget,
radically different map: \textbf{the allocator is path-dependent.} Start-low-and-upgrade is a
distinct strategy from start-at-floor-and-downgrade, and it has not been run.

%% ============================================================ 10
\section{Verdicts across all three panels}

DynQuant on the left. Shipped-convention baselines, so DynQuant is the smaller arm in every row.

\subsection{4.25 bits (3.77\x)}

\begin{center}\small
\begin{tabular}{@{}lrrrl@{}}
\toprule
vs & Qwen/CaseHOLD & Qwen/Banking77 & Mistral/Banking77 & replicated? \\
\midrule
GPTQ 4b & $-0.49$ (0.10) & $-0.42$ (0.14) & $-0.06$ (0.85) & yes --- never separated \\
AWQ 4b & $+0.08$ (0.85) & $+0.16$ (0.64) & $-0.23$ (0.21) & yes --- never separated \\
RTN 4b & $+0.32$ (0.34) & \best{$+1.14$ (\e{7.3}{-4})} & $-0.06$ (0.83) & \textbf{no} --- one panel only \\
NF4 & $+0.24$ (0.45) & $-0.23$ (0.48) & $-0.06$ (0.85) & yes --- never separated \\
bf16 & $-0.49$ (0.099) & \best{$-0.75$ (\e{2.7}{-3})} & $-0.19$ (0.26) & \textbf{no} --- one panel only \\
\bottomrule
\end{tabular}
\end{center}

\subsection{3.25 bits (4.92\x)}

\begin{center}\small
\begin{tabular}{@{}lrrrl@{}}
\toprule
vs & Qwen/CaseHOLD & Qwen/Banking77 & Mistral/Banking77 & replicated? \\
\midrule
GPTQ 3b & \best{$-2.28$ (\e{7.0}{-10})} & \best{$-1.04$ (0.011)} & \best{$-0.78$ (\e{7.1}{-3})} & \textbf{yes --- GPTQ ahead on all three} \\
AWQ 3b & \best{$+3.33$ (\e{7.5}{-14})} & \best{$+2.47$ (\e{3.3}{-7})} & $-0.42$ (0.14) & partly --- two of three \\
RTN 3b & \best{$+21.17$} & \best{$+29.48$} & $+0.26$ (0.42) $\dagger$ & direction on all three \\
bf16 & \best{$-3.05$ (\e{4.9}{-17})} & \best{$-2.31$ (\e{6.8}{-10})} & \best{$-0.78$ (\e{3.2}{-3})} & yes --- separated on all three \\
\bottomrule
\end{tabular}
\end{center}

$\dagger$ The Mistral cell compares a 3.25-bit DynQuant against a 3.62-bit RTN, so DynQuant is
giving away 10\pct{} of the size. At matched bytes it separates: \best{$+0.71$}, \e{7.2}{-3}.

\subsection{What replicated, and what did not}
\label{sec:replication}

\subsubsection*{Replicated across all three panels}

\begin{enumerate}[leftmargin=1.4em]
\item \textbf{GPTQ wins the 3-bit tier.} $-2.28$ / $-1.04$ / $-0.78$, separated every time. The
single most robust conclusion in the campaign, and it is the one unfavourable to DynQuant.
\textbf{Superseded on Qwen/CaseHOLD by a later rebuild} --- sensitivity allocation, an
$\mathbb{E}[x^2]$-weighted clip objective, a row-partitioned tied embedding and per-row body
widths turn that $-2.28$ into \best{$+1.54$ at 4.5\pct{} fewer bytes}, $p < 0.0001$. The other two
panels have not been re-run and stand as written.

\item \textbf{The 4-bit tier is a tie.} Across three panels and four baselines, 12 comparisons at
4-bit-class budgets produce exactly one separation (RTN on Qwen/Banking77). DynQuant achieves that
tie at 1.73\x{} / 1.73\x{} / 1.08\x{} fewer bytes.

\item \textbf{Quantization to 3.25 bits costs real accuracy.} Separated from bf16 on all three
panels. No arm of DynQuant below 4 bits is lossless.

\item \textbf{DynQuant beats RTN at 3 bits} --- and this one needed matched bytes to see. On the
two Qwen panels the shipped-convention comparison is overwhelming ($+21.17$, $+29.48$). On Mistral
it \emph{fails} ($+0.26$, $p = 0.42$) --- but that comparison gives away 10\pct{} of the size. At
matched bytes it separates: \best{$+0.71$}, \e{7.2}{-3}. The direction holds on all three panels;
only the magnitude collapses, because RTN does not collapse on Mistral.
\end{enumerate}

\subsubsection*{Replicated on both Qwen panels, absent on Mistral}

\begin{enumerate}[leftmargin=1.4em,start=5]
\item \textbf{DynQuant beats AWQ at 3 bits} ($+3.33$, $+2.47$) at roughly half the bytes --- and
$+3.39$ against AWQ's matched-byte \code{\_head} arm on CaseHOLD, so the win is not an artefact of
the size mismatch. On Mistral at matched bytes the two are a dead tie: $+0.03$, 29/28 flips,
$p = 1.00$. This one genuinely does not carry over.
\end{enumerate}

The pattern in both cases is the same and it is not a DynQuant property: \textbf{RTN and AWQ do
not collapse on Mistral}, so there is no collapse for DynQuant to avoid. RTN at 3 bits loses 1.04
points there against 24--32 on Qwen. The Mistral panel is not underpowered --- it resolves that
1.04 at $p = $ \e{5.8}{-5}, and it resolves DynQuant's $+0.71$ over RTN at $p = $ \e{7.2}{-3} ---
it is \emph{compressed}: the entire span from best method to worst is 1.08 points, so a 25-point
advantage and a 0.7-point advantage cannot look different in kind. \textbf{It lacks dynamic
range, not sample size.}

\subsubsection*{Did not replicate across tasks on the same model}

\begin{enumerate}[leftmargin=1.4em,start=6]
\item \textbf{DynQuant 4.25 vs RTN 4b} separated on Banking77 ($+1.14$) and not on CaseHOLD
($+0.32$). The mechanism is worth stating precisely, because it is not the flattering one:
DynQuant's 4.25-bit arm did not get \emph{better} on Banking77 --- it got worse, by 0.26 points.
It separated from RTN there because RTN got worse \emph{faster}, losing 1.07 where DynQuant lost
0.26. \textbf{The gap opens when the task punishes uniform rounding, not when the signal finds
something extra.}
\end{enumerate}

\subsubsection*{One negative result worth recording as a result}

\textbf{The bit map is mostly a property of the architecture, not the task.} At 4.25 bits the
CaseHOLD and Banking77 maps agree on 177 of 187 widths; 86--89\pct{} of off-uniform decisions are
task-invariant. A method sold on task-adaptive allocation produces a nearly task-invariant map at
moderate budgets.

%% ============================================================ 11
\section{Defects found and fixed during this campaign}
\label{sec:defects}

Five, all found by guards or by measurement rather than by reading. Recorded because the fixes are
now the harness and because three of them would each have produced a wrong published number.

\subsection*{1. The tied-embedding accounting defect}

\code{ignore=["lm\_head"]} makes a ``4-bit'' baseline measure 7.3605 bits on a tied model. Found
by computing bytes independently instead of trusting the label. Fixed by adding the seven
\code{\_head} arms so every claim has a matched-byte version. \emph{Consequence if missed: every
byte-ratio claim in the report would have been inflated by 1.7\x.}

\subsection*{2. A stale bit map on the Mistral panel}

Both Mistral DynQuant arms were originally quantized through a map built from a
\textbf{superseded} fine-tune's statistics. The fine-tune had been re-run; the map had not. Every
skip-if-output-exists resume guard stayed true because every output still existed ---
\textbf{existence cannot detect staleness, ordering can}. Repaired by rebuilding the map and
re-running both arms.

\emph{The defect turned out to be empty in consequence, and this is measured, not assumed:} at
4.25 bits the corrected map differs from the stale one on \textbf{0 of 226 widths}, and the re-run
arm is \textbf{bit-identical} --- $+0.00$, 0/0 flips, $p = 1.00$, across 3{,}080 examples. At 3.25
bits the maps differ on 24 of 226 widths and the accuracy difference is \textbf{$+0.03$, CI
$[-0.31, +0.38]$, 15/14 flips, $p = 1.00$}. Both stale rows are retained in Panel C as
\code{\_ft1map} so the claim is checkable.

\subsection*{3. Records written into another model's directory}

\code{RUN\_DIR} is derived from \code{DQ\_MODEL} and \code{DQ\_TASK}, not from \code{--model}. A
hand-written driver pinned the run directory in a shell variable and passed \code{--model}
explicitly --- pinning only the \emph{input}. With \code{DQ\_MODEL} unset, four Mistral arms wrote
into the \textbf{Qwen} directory, and the same defaulted \code{MODEL\_ID} handed a \textbf{Qwen
tokenizer to Mistral weights}, producing out-of-vocabulary token ids and a \code{device-side
assert} on every arm --- \emph{after} each 7B quantize had already been paid for. Three files were
written under Qwen names, two of them overwriting real Qwen records. All accuracy records survived,
because \code{record()} runs after eval and every one of these arms crashed \emph{during} eval.

The misplaced files are quarantined with a README, deliberately \textbf{not} moved into the
directory they were meant for: a \code{\_quant.json} there with no matching accuracy record reads
as a partially-completed arm, which is the same stale-artifact trap that produced defect 2. The
lesson, generalized: \textbf{passing the input explicitly does not pin the output.}

\subsection*{4. A guard that fired for the wrong reason}

The first repair attempt aborted with ``the map was never rebuilt'', inferred from the rebuilt map
being byte-identical to the old one. It was wrong --- the map \emph{had} been rebuilt, and at 4.25
bits rebuilding it correctly produces the same file. \textbf{A content diff cannot separate ``never
rebuilt'' from ``rebuilt identically''; only mtimes can.} This is the second time in this project
that a guard's stated reason for firing turned out to be a hypothesis rather than a diagnosis, and
the general form is worth keeping: \textbf{a guard reports that its condition tripped, never why.}

\subsection*{5. A silent container restart}

The repair chain died at 12:58:19 with no traceback, no OOM record and the GPU freed.
\code{memory.events} showed \code{oom\_kill 0} and a peak of 4.15 GB against a 208 GB limit; no XID
or ECC errors. The cause was the host restarting the container at 13:00:31 --- \code{.log.old}
files from 08:32 show it had happened earlier the same day. It was invisible for 44 minutes because
the log still \emph{said} \code{992/3080} while the box had come back underneath it.
\textbf{Silence is not success}; liveness is now checked by log age plus process presence, not by
content.

\subsection{The guards these bought}

Defects 2 and 3 are now checked inside the quantize stage rather than in a runner, so they apply to
every arm regardless of which driver launches it. To be exact about the timeline: \textbf{the arms
in this report were not produced under these guards.} They were produced under a runner-level
\code{assert\_rundir.py} check plus a hand-rebuilt map, and the in-script guards were written
afterwards so the \emph{next} campaign cannot repeat either defect. What protects the numbers here
is the measurement in defect 2 above, not the guard.

\begin{itemize}[leftmargin=1.4em]
\item \textbf{Output-directory check} --- \code{--model} must resolve inside the \code{RUN\_DIR}
that \code{common} itself computed. Catches defect 3 before a single weight is loaded.
\item \textbf{Map-is-youngest check} --- the bit map must be newer than both the weights and the
stats it describes, with \code{newest\_mtime} reaching inside directories because a directory's own
mtime does not move when a file inside it is overwritten in place. Catches defect 2.
\item \code{--allow-stale} downgrades both to warnings, for the one legitimate case: deliberately
measuring a stale pairing to quantify what the staleness cost. That is how the \code{\_ft1map} rows
exist.
\end{itemize}

One trap encountered while writing the freshness check is recorded in the source, because it would
have been a plausible-looking guard that broke every clean run: an earlier draft also asserted
$\mathrm{mtime(stats)} \geq \mathrm{mtime(weights)}$. But stats are written \emph{throughout}
training and the weights are saved at the \emph{end}, so stats-older-than-weights is the normal
case --- on the Mistral run, by eleven seconds. Only ``the map is the youngest artifact'' is a
valid assertion.

\subsection{The determinism check on the two clobbered records}

The two overwritten \code{\_quant.json} files were regenerated by re-running both Qwen/Banking77
DynQuant arms end to end, which doubles as a determinism test: the accuracy records were never
damaged, so the re-run has a published number to reproduce. \textbf{Both reproduce exactly.}

\begin{center}\small
\begin{tabular}{@{}lrrrlrr@{}}
\toprule
arm & published & re-run & delta & CI95 & flips & $p$ \\
\midrule
dq\_4p25 & 92.6623\pct & 92.6623\pct & $+0.00$ & $[+0.00, +0.00]$ & 0/0 & 1.00 \\
dq\_3p25 & 91.1039\pct & 91.1039\pct & $+0.00$ & $[+0.00, +0.00]$ & 0/0 & 1.00 \\
\bottomrule
\end{tabular}
\end{center}

Not merely equal accuracy --- \textbf{0/0 flips across 3{,}080 examples each}, so every individual
prediction is identical. Allocation, MSE clipping and evaluation are all reproducible from the
seed, and the numbers in Panel B were not a lucky draw.

\subsection{The guards are tested, and the test found a hole in them}

\code{test\_guards.sh} exercises all four paths. Every case is designed to exit \emph{before}
\code{load\_model}, so validating a guard costs seconds rather than a 7B quantize.

\begin{center}\small
\begin{tabular}{@{}lll@{}}
\toprule
case & expected & result \\
\midrule
map backdated 1 h before the weights & abort rc=3, no model load & \textbf{pass} --- rc=3, no load line \\
\code{DQ\_MODEL} unset, \code{--model} at Mistral & abort rc=3 naming the Qwen dir & \textbf{pass} \\
correct map + correct \code{DQ\_MODEL} & \code{provenance: ok}, proceed & \textbf{pass} (rc=2 on bogus target) \\
stale map + \code{--allow-stale} & warn, proceed & \textbf{pass} \\
\bottomrule
\end{tabular}
\end{center}

6 of 6 assertions pass and no test writes a record. Note the timeline honestly: \textbf{this suite
ran after the campaign}, so it certifies the guards, not the arms.

\textbf{The first version of the test failed, and its failure is the more useful result.} It used
the retained \code{stage4\_bitmaps.ft1.json} as the stale map --- the actual stale artifact from
defect 2. The guard passed it, and the run went on to quantize and evaluate for real before it was
killed. The guard was right and the test was wrong: that file is a \emph{copy}, made on 2026-07-30
11:50, so its mtime is a day \textbf{younger} than the 2026-07-29 03:08 weights it misdescribes.

\begin{aside}
\textbf{An mtime guard cannot see through a copy, a \code{git checkout}, an \code{rsync} without
\code{-t}, a container image rebuild, or a download.} Each of those refreshes the mtime while
preserving stale content. It catches the failure mode that actually occurred --- a map left
untouched while its inputs were regenerated --- and it is blind to a stale map that has since been
moved.
\end{aside}

Closing that would need the map to record a content hash of the stats it was built from, and the
check to compare hashes rather than times. That is the correct fix and it is not implemented. Until
it is, the guard is a tripwire for one specific accident, not a provenance system --- and defect 4
already showed that a content diff alone cannot substitute, because it cannot separate ``never
rebuilt'' from ``rebuilt identically''. Times and hashes each catch what the other misses; the
guard currently has only one of them.

%% ============================================================ 12
\section{Limitations}
\label{sec:limits}

Stated plainly, because most of them bound the conclusions above more tightly than the $p$-values
do.

\begin{enumerate}[leftmargin=1.4em]
\item \textbf{Two models.} One 1.88\,B tied and one 7.25\,B untied. Every claim that differs
between them is confounded by scale \emph{and} by the tie, and no arm separates the two.

\item \textbf{Two tasks, both classification.} Accuracy on a fixed answer set is a coarse metric.
No generative task, no perplexity, no long-context measurement. A method could preserve
classification accuracy and still damage generation.

\item \textbf{One clean win, on one tier.} The 6.63-bit result is a single (model, task, budget)
point. It has not been reproduced on Banking77 or on Mistral, and the equal-accuracy arm at 7.36
bits suggests the win sits at the edge of where degradation begins rather than in the middle of a
trend.

\item \textbf{The runtime evidence covers DynQuant only, and its decode numbers are not tightly
measured.} \S\ref{sec:packed} confirms the bytes, the VRAM and the kernel-vs-simulated equivalence
on Mistral, so the accuracy comparisons are not undermined by the packed path. What it does not do
is measure any \emph{baseline} packed: GPTQ, AWQ, RTN and NF4 were scored through
\code{llm-compressor}'s own simulated output and never run through a serving kernel here, so
\textbf{no speed or VRAM comparison against a baseline exists in this report} --- the 0.87\x{}
decode figure is DynQuant against bf16, not DynQuant against GPTQ, and Marlin-class baseline
kernels would very likely beat it.

\item \textbf{The ablation is one model, one task, two budgets.} The 88/12 split is measured at
3.25 bits on Qwen/CaseHOLD. It is not known whether the signal's share grows at 3 bits or below,
and the path-dependence lever suggests the allocator's own share is not even fixed.\footnote{Phase
2 answered part of this: at \emph{row} granularity on the same panel, the signal's shuffled control
loses 1.28 points and the signal itself is worth $+2.31$ --- so the 12\pct{} figure is scoped to
module granularity, not to the method.}

\item \textbf{Baselines are one library's implementation at one configuration.} 256 calibration
rows, \code{group\_size} 128, \code{dampening\_frac} 0.01. A better-tuned GPTQ would widen its lead
at 3 bits; a more aggressive AWQ configuration might close its gap. No baseline was tuned.

\item \textbf{The Mistral panel is compressed, not powerful.} Its whole best-to-worst span is 1.08
points. It does separate six comparisons, all in the 3-bit tier, so it is not underpowered --- but
no 4-bit-class row there can distinguish a large effect from a negligible one, and those rows
should not be read as evidence either way.
\end{enumerate}

%% ============================================================ 13
\section{What this means for the method}

The comparison is not flattering everywhere, and the shape it does have is specific.

\textbf{DynQuant's real contribution, as measured, is the allocation prior --- not the training
signal.} 88\pct{} of its 25.78-point margin over byte-matched RTN comes from role floors,
structural protection and honest byte accounting; those are architectural facts that require no
training run to obtain. The signal contributes a separated but small 3.16 points, and only where
the budget is tight enough to force many decisions. A user who wants most of DynQuant's benefit
does not need to instrument a fine-tune.

\textbf{The 3-bit tier is the method's frontier and it currently loses there.} GPTQ is smaller and
better at 3 bits on all three panels. The mechanism is legible: DynQuant chooses widths well and
then rounds naively, while GPTQ chooses nothing and compensates error. \textbf{These are
orthogonal.} Nothing about signal-driven width allocation precludes a better encoder, a finer
allocation partition, or a sensitivity model in units of loss --- and the obvious next experiment is
to stop treating the 3-bit gap as structural.

\begin{aside}
\textbf{What Phase 2 did with this paragraph.} It took the last sentence literally and rebuilt the
three components DynQuant owns --- the sensitivity model, the allocation partition and the clip
objective --- while deliberately \emph{not} adopting inverse-Hessian error compensation, which
would have required a calibration set and would have made the method a reimplementation of GPTQ.
The result on this panel is $+1.54$ at 4.5\pct{} fewer bytes. See the companion report.
\end{aside}

\textbf{The compression is real and the speed is the next problem, in that order.} The packed
checkpoints match the accounted bytes exactly, cut peak VRAM 3.55--4.55\x, and reproduce the
simulated accuracies to within a handful of predictions --- which retires the paper's own concession
that the savings are storage-only, and retires the worry that the simulated evaluation was measuring
something the kernels do not do. What replaces it is a narrower complaint: decode is 10--15\pct{}
slower than bf16 at batches 1--8 and falls apart at 32. Both causes are diagnosed rather than
guessed --- a fixed ${\sim}15$ \textmu s per-module launch tax that CUDA Graphs removes, and a
\code{GEMV\_MAX\_ROWS} of 8 that sends batch 32 down a dequantize-then-\code{F.linear} fallback
reading more memory than bf16. Neither is a property of signal-driven bit allocation. But the honest
position today is that DynQuant trades throughput for footprint, and no baseline was benchmarked
packed, so its speed against a Marlin-class GPTQ kernel is unknown and should not be assumed
favourable.

\textbf{The clean claim the method could make at the close of Phase 1} is the 6.63-bit one: lossless
at 2.42\x{} compression where GPTQ at identical bytes is not. That is a narrower claim than the
paper's headline and it is the one that survives a matched-byte comparison.

%% ============================================================ 14
\section{Reproduction}

Everything runs from \code{experiments/four\_point/} on a single 80 GB GPU.

{\footnotesize
\begin{verbatim}
# 1. environment for the baselines (llm-compressor + bitsandbytes)
python -m venv venv-cmp && venv-cmp/bin/pip install llmcompressor bitsandbytes

# 2. pin BOTH variables -- RUN_DIR is derived from them, not from --model
export DQ_MODEL=Qwen/Qwen3.5-2B-Base
export DQ_TASK=casehold
RUN=/workspace/runs/qwen3_5_2b_base_casehold

# 3. verify python and shell agree on the output directory before anything runs
python assert_rundir.py "$RUN"

# 4. baselines, shipped convention and matched-byte
bash run_baselines.sh          # gptq/awq/rtn at 4 and 3 bits
bash run_tiedhead.sh           # the same with ignore=[]  -> the _head arms
python stage8_bnb.py --name stage8_nf4_4b

# 5. DynQuant arms
python stage4_allocate.py --targets 4.25,3.25
python stage5_quantize.py --target 4.25 --name stage8_dq_4p25
bash run_isosize.sh            # the equal-byte arms at 7.36 and 6.63 bits

# 6. the signal ablation
python make_control_map.py     # uniform scores, everything else fixed
python stage5_quantize.py --bitmaps "$RUN/stage4_bitmaps.control.json" \
                          --target 3.25 --name stage8_dq_ctl3p25

# 7. read the results
python dump_all.py                                   # every arm, every panel
python pairs.py "$RUN" dq_3p25:gptq_3b_head          # McNemar on any pair
\end{verbatim}}

\code{pairs.py} takes the run directory \textbf{first} and prefixes \code{stage8\_} itself, so pass
bare arm names.

\subsection{File map}

\begin{center}\small
\begin{tabular}{@{}ll@{}}
\toprule
file & what it is \\
\midrule
\code{stage4\_allocate.py} & the knapsack; runs twice per target (real scores + uniform control) \\
\code{stage5\_quantize.py} & applies a map and scores it; holds both provenance guards \\
\code{stage8\_baselines.py} & llm-compressor GPTQ/AWQ/RTN, with the rounding verification \\
\code{stage8\_bnb.py} & bitsandbytes NF4 \\
\code{test\_guards.sh} & validates both provenance guards; every case exits before model load \\
\code{bitmap\_diff.py} & width-level diff between two maps \\
\code{baselines\_table.py} & regenerates the panel tables from the records \\
\code{RESULTS-external-comparison.md} & the working document, with derivations and failed attempts \\
\code{RESULTS-mistral7b-banking77.md} & the Mistral campaign, including the packed VRAM/speed data \\
\code{RESULTS.md} & the Qwen campaign: kernel parity, sanitizer sweep, bandwidth gates \\
\bottomrule
\end{tabular}
\end{center}

\vspace{1.5em}
\begin{center}
{\color{dqblue}\rule{0.5\textwidth}{0.8pt}}
\end{center}
\vspace{0.5em}

{\small\itshape
All 45 comparison arms are measured and every verification job has landed. The last one --- a
re-run of the two Qwen/Banking77 arms whose \code{\_quant.json} was overwritten by defect 3 ---
reproduced both published numbers with 0/0 prediction flips. \S\ref{sec:packed} was added last,
after diffing its bit maps against the shipped ones, and it replaced an earlier limitation that
wrongly claimed no runtime measurement existed. Nothing in this report is pending.\par}

\end{document}
